\begin{tikzpicture}[font=\small]
  % ---------- periodic DT signal ----------
  \begin{scope}
    \draw[ax] (-4.6,0) -- (4.4,0) node[below left]{$n$};
    \draw[ax] (0,-0.2) -- (0,1.8);
    \node[note, anchor=south east] at (-0.1,1.45){$x[n]$};
    \foreach \k in {-2,-1,0,1}{
      \foreach \n/\h in {0/1.3, 1/0.75, 2/0.35, 3/0.75}{
        \stemat{ssblue}{(\k*4+\n)*0.5}{\h}
      }
    }
    \draw[helper] (2.0,-0.25) -- (2.0,1.6);
    \draw[helper] (0,-0.25) -- (0,1.6);
    \draw[{Stealth[length=4pt]}-{Stealth[length=4pt]}, ssred] (0,-0.42) -- (2.0,-0.42);
    \node[ssred, font=\footnotesize, anchor=north] at (1.0,-0.48){$N=4$};
    \node[note, anchor=north, align=center] at (0,-1.05){周期序列，周期 $N$};
  \end{scope}

  % ---------- its DTFS coefficients, also periodic ----------
  \begin{scope}[xshift=10.4cm]
    \draw[ax] (-4.6,0) -- (4.4,0) node[below left]{$k$};
    \draw[ax] (0,-0.2) -- (0,1.8);
    \node[note, anchor=south east] at (-0.1,1.45){$a_k$};
    \foreach \k in {-2,-1,0,1}{
      \foreach \n/\h in {0/0.79, 1/0.48, 2/0.29, 3/0.48}{
        \stemat{ssteal}{(\k*4+\n)*0.5}{\h}
      }
    }
    \draw[helper] (2.0,-0.25) -- (2.0,1.6);
    \draw[helper] (0,-0.25) -- (0,1.6);
    \draw[{Stealth[length=4pt]}-{Stealth[length=4pt]}, ssred] (0,-0.42) -- (2.0,-0.42);
    \node[ssred, font=\footnotesize, anchor=north] at (1.0,-0.48){$N=4$};
    \node[note, anchor=north, align=center] at (0,-1.05){系数序列也是周期 $N$ 的};
  \end{scope}

  \node[note, anchor=north, align=center] at (5.2,-2.4)
    {$a_k=\dfrac{1}{N}\displaystyle\sum_{n=\langle N\rangle}x[n]e^{-jk(2\pi/N)n}$,\qquad
     $x[n]=\displaystyle\sum_{k=\langle N\rangle}a_k e^{jk(2\pi/N)n}$};

  \node[note, anchor=north, align=center] at (5.2,-3.9)
    {与 CTFS 的决定性差别：谐波只有 \textbf{$N$ 个不同的}（$e^{j(k+N)(2\pi/N)n}=e^{jk(2\pi/N)n}$），
     所以求和是\textbf{有限}的，\\ 也就没有收敛性问题、没有吉布斯现象 —— DTFS 是这四种傅里叶表示里唯一「完全没有分析学麻烦」的一个};
\end{tikzpicture}
